\documentclass[11pt]{article}
\usepackage{amsmath,amssymb,amsthm}
\usepackage[margin=1.1in]{geometry}
\usepackage{hyperref}

\newtheorem{proposition}{Proposition}
\newcommand{\sigmoid}{\sigma}
\newcommand{\xbar}{\bar{x}}

\title{Bernoulli Maximum-Likelihood Estimation:\\
Closed Form, Gradient Descent, SGD, and Autograd}
\author{Yuval Lavie}
\date{\today}

\begin{document}
\maketitle

\begin{abstract}
We estimate the parameter of a Bernoulli distribution from i.i.d.\ samples
four ways: the theoretical (closed-form) maximum-likelihood solution;
gradient descent on the negative log-likelihood with a gradient derived by
hand; stochastic gradient descent, which replaces that gradient with a
one-sample unbiased estimate; and the same SGD with the gradient computed
by PyTorch's autograd. The companion script \texttt{mle.py} runs all four
and verifies, by driving the two SGD variants through the identical sample
schedule, that autograd reproduces the hand-derived gradient to
floating-point precision.
\end{abstract}

\section{Setup: from the likelihood to the NLL, step by step}

Let $x_1,\dots,x_n \overset{\text{iid}}{\sim} \mathrm{Bernoulli}(p)$ with
$x_i \in \{0,1\}$ and unknown $p \in (0,1)$.

\paragraph{Step 1: the likelihood is a product of two kinds of factors.}
By independence, the probability of observing exactly this sample is
\begin{equation}
  L(p) \;=\; \prod_{i=1}^{n} \Pr(x_i)
       \;=\; \prod_{i=1}^{n} p^{x_i}(1-p)^{1-x_i}.
\end{equation}
Each factor is just a case split written as one expression: it equals $p$
when $x_i = 1$ and $1-p$ when $x_i = 0$. So if the sample contains
$m = \sum_i x_i$ ones and $n - m$ zeros, collecting the factors gives
\begin{equation}
  L(p) \;=\; p^{\,m}\,(1-p)^{\,n-m}.
  \label{eq:likelihood}
\end{equation}
Already here the data enters only through the count $m$: how the ones and
zeros are ordered is irrelevant. ($m$, equivalently $\xbar = m/n$, is a
\emph{sufficient statistic} for $p$.)

\paragraph{Step 2: the logarithm turns the product into a sum.}
$\log$ is strictly increasing, so it preserves the location of a maximum:
$\arg\max_p L(p) = \arg\max_p \log L(p)$. Applying it to
\eqref{eq:likelihood},
\begin{equation}
  \ell(p) \;=\; \log L(p) \;=\; m\log p + (n-m)\log(1-p).
  \label{eq:loglik}
\end{equation}
This is the whole reason for working in logs: powers become products,
products become sums, and sums are easy to differentiate (and do not
underflow numerically the way a product of $5000$ numbers in $(0,1)$
would).

\paragraph{Step 3: normalize and negate --- neither moves the optimum.}
Two cosmetic transformations bring \eqref{eq:loglik} to the standard
machine-learning form. Dividing by the constant $n > 0$ rescales the
function but does not move its argmax. Negating flips maximization into
minimization (the ML convention is to \emph{minimize losses}). With
$\xbar = m/n$, the average negative log-likelihood (NLL) is
\begin{equation}
  \mathcal{L}(p) \;=\; -\tfrac1n\,\ell(p)
  \;=\; -\xbar\log p - (1-\xbar)\log(1-p).
  \label{eq:nll}
\end{equation}
Chaining the three steps:
\begin{equation}
  \arg\max_p L(p)
  \;=\; \arg\max_p \ell(p)
  \;=\; \arg\min_p \mathcal{L}(p),
  \label{eq:equiv}
\end{equation}
so maximizing the likelihood and minimizing the NLL are the same problem,
each step preserving the optimizer.

\section{Route 1: the theoretical solution}

The cleanest derivation works directly on the log-likelihood
\eqref{eq:loglik}, where the counts keep the algebra honest.

\begin{proposition}
Assume $0 < m < n$. Then $\hat{p}_{\mathrm{MLE}} = m/n = \xbar$, and it is
the unique maximizer of $\ell$ on $(0,1)$.
\end{proposition}

\begin{proof}
Differentiate \eqref{eq:loglik} term by term:
\begin{equation*}
  \ell'(p) \;=\; \frac{m}{p} - \frac{n-m}{1-p}.
\end{equation*}
Set it to zero and clear denominators:
\begin{equation*}
  \frac{m}{p} = \frac{n-m}{1-p}
  \;\;\Longrightarrow\;\;
  m(1-p) = (n-m)\,p
  \;\;\Longrightarrow\;\;
  m - mp = np - mp
  \;\;\Longrightarrow\;\;
  m = np,
\end{equation*}
hence
\begin{equation}
  \boxed{\;\hat{p}_{\mathrm{MLE}} \;=\; \frac{m}{n} \;=\; \xbar\;}
  \label{eq:phat}
\end{equation}
--- the fraction of ones in the sample. Uniqueness:
$\ell''(p) = -m/p^{2} - (n-m)/(1-p)^{2} < 0$, so $\ell$ is strictly concave
on $(0,1)$ and the stationary point is the global maximum.
\end{proof}

(If $m = 0$ or $m = n$, every factor of the likelihood pulls in the same
direction: $L$ is monotone on $(0,1)$ and its supremum sits at the
boundary, which is again $m/n$ --- so the formula $\hat p = m/n$ survives
the edge cases.) Since the three optimization problems of Section~1 share
their optimizer, $\xbar$ also minimizes the NLL \eqref{eq:nll} --- the
form the gradient routes below will descend.

\section{Route 2: gradient descent}

\paragraph{One coordinate change --- not a new objective.}
Gradient methods want an unconstrained variable, but $p$ lives in $(0,1)$.
So we evaluate the \emph{same} NLL \eqref{eq:nll} at
\begin{equation}
  p = \sigmoid(t) = \frac{1}{1+e^{-t}}, \qquad t \in \mathbb{R},
  \qquad
  F(t) \;=\; \mathcal{L}\bigl(\sigmoid(t)\bigr).
  \label{eq:reparam}
\end{equation}
Nothing about the target function changes; we have only composed it with a
smooth bijection of the search space. (By the \emph{invariance} of maximum
likelihood, the optimum maps back: $\hat p = \sigmoid(\hat t\,)$.)

\paragraph{The gradient, by the chain rule.}
Using $\sigmoid'(t) = \sigmoid(t)\bigl(1-\sigmoid(t)\bigr)$ and writing
$s = \sigmoid(t)$ for the moment,
\begin{equation*}
  F'(t)
  \;=\; \mathcal{L}'(s)\,\sigmoid'(t)
  \;=\; \left[-\frac{\xbar}{s} + \frac{1-\xbar}{1-s}\right] s(1-s)
  \;=\; -\xbar(1-s) + (1-\xbar)\,s
  \;=\; s - \xbar ,
\end{equation*}
hence
\begin{equation}
  F'(t) \;=\; \sigmoid(t) - \xbar .
  \label{eq:grad}
\end{equation}
Every term cancelled except the residual: current estimate minus sample
mean. Moreover $F''(t) = \sigmoid'(t) > 0$, so $F$ is convex in $t$ ---
one basin.

\paragraph{Gradient descent.}
GD iterates the full gradient, every step:
\begin{equation}
  t \;\leftarrow\; t - \eta\,\bigl(\sigmoid(t) - \xbar\bigr).
  \label{eq:gd}
\end{equation}
This is deterministic: with a sane step size it converges monotonically to
the unique fixed point $\sigmoid(t^{\ast}) = \xbar$ --- the closed form of
Route~1, now reached by iteration. For this toy problem the ``full
gradient'' is free ($\xbar$ is precomputed); in real models it costs a
pass over the entire dataset per step, which is what motivates the next
route.

\section{Route 3: stochastic gradient descent}

SGD replaces the full gradient \eqref{eq:grad} with a \emph{one-sample
estimate}: draw an index $i$ and use
\begin{equation}
  g_i(t) \;=\; \sigmoid(t) - x_i .
  \label{eq:sgdgrad}
\end{equation}
It is unbiased --- averaging over a uniformly drawn $i$,
$\mathbb{E}_i\!\left[g_i(t)\right] = \sigmoid(t) - \xbar = F'(t)$ --- so
each cheap step points in the right direction \emph{on average}. The
script sweeps reshuffled passes over the data, updating
$t \leftarrow t - \eta\, g_i(t)$ one sample at a time.

\paragraph{The price of stochasticity, and Polyak--Ruppert averaging.}
With a constant step size the SGD iterate never converges: it reaches a
stationary \emph{distribution}, a noise ball around the optimum whose
radius grows with $\eta$ and the per-sample variance. The estimator is
therefore taken as the average of the final pass's iterates,
$\hat{p}_{\mathrm{SGD}} = \tfrac1K \sum_k \sigmoid(t_k)$, which cancels
the hovering and lands within a few $10^{-3}$ of $\xbar$.

\paragraph{Mini-batches (later).}
Averaging $g_i$ over a small batch $B$ gives $\sigmoid(t)-\xbar_B$: the
same expectation with variance reduced by $|B|$. Everything in this note
goes through unchanged; batching is an engineering trade between noise and
cost per step, not new mathematics.

\section{Route 4: SGD with autograd}

PyTorch's \texttt{binary\_cross\_entropy\_with\_logits}$(t, x_i)$
evaluates exactly the one-sample NLL at $p = \sigmoid(t)$, i.e.\ the
composition $F_i(t) = \mathcal{L}_i(\sigmoid(t))$ of Route 2 restricted to
sample $i$. Calling \texttt{loss.backward()} applies reverse-mode
automatic differentiation --- a mechanized chain rule --- so it must
produce \texttt{t.grad} $= \sigmoid(t) - x_i$, identically
\eqref{eq:sgdgrad}. \texttt{torch.optim.SGD} then performs the same update
$t \leftarrow t - \eta\,\texttt{t.grad}$. Autograd changes who does the
calculus, not the mathematics.

\paragraph{Proof by computation.}
The script drives Routes 3 and 4 through the \emph{identical} sample
schedule, the same $\eta$, and the same initialization $t_0 = 0$ (in
\texttt{float64} on both sides). Under those conditions the two
trajectories can differ only if the gradients differ; the script asserts
$\max_k |\sigmoid(t_k^{\text{hand}}) - \sigmoid(t_k^{\text{autograd}})| <
10^{-10}$, which passes --- an end-to-end check that
\eqref{eq:sgdgrad} is exactly what autograd computes.

\section{Experiment}

With $p = 0.3$ and $n = 5000$ (seed 7): the closed form gives
$\hat{p} = m/n \approx 0.2978$; GD converges to it deterministically (the
gap shrinks geometrically); SGD, Polyak-averaged over the last pass, lands
within a few $10^{-3}$; and the two SGD trajectories coincide to
$10^{-10}$. The figure produced by \texttt{mle.py} shows, left to right:
the strictly convex NLL landscape; GD's smooth convergence to $m/n$; and
SGD's noise ball around it, with the autograd trajectory overlaid and
indistinguishable.

Four routes, one destination: the likelihood surface --- not the
optimizer --- decides where estimation lands.

\end{document}
